\section{Discussion}

\subsection{Principal Findings}

The study yields five main findings. First, a multimodal Coswara model can
separate labels well in a participant-disjoint internal test, but that estimate
does not determine how a cough model will behave in another collection system.
Second, the external decline persisted across four conventional classifiers,
WavLM, and CNN--BiGRU. The consistency across representation and model families
makes it less plausible that the transfer gap is explained only by one weak
classifier. Third, feature selection was highly unstable across calendar
periods. Fourth, recording context and symptoms were strongly associated with
the label, and shuffled-label retraining removed that association. Fifth, on
the small aligned cohort, adding audio to metadata did not yield a precise or
statistically supported AUROC improvement.

These results do not imply that every internal result is invalid. They show
that the relevant claim must match the evaluation. AUROC 0.897 is evidence of
internal discrimination for one selected multimodal protocol. It is not an
estimate of multimodal performance on COUGHVID, because COUGHVID has no matched
breathing and speech inputs. Conversely, external AUROC near 0.5 is evidence of
poor portability under this source-target pair; it is not proof that no
disease-related acoustic information exists in any population.

\subsection{Relation to Prior Evidence}

The findings are consistent with the realistic-evaluation studies of Han
\emph{et al.} and Coppock \emph{et al.}, which showed that participant design,
demographic balance, symptoms, and enrolment change measured audio performance
\cite{han2022realistic,coppock2024audio}. Our study extends that concern within
a public multimodal Coswara workflow by adding forward and reverse calendar
tests, temporal feature-set overlap, explicit source-to-COUGHVID transfer,
deep-representation transfer, calibration, and target operating-point
analysis.

The comparison with the HST and DNDT/DNDF papers requires protocol discipline.
HST reported strong 10-fold results for dataset-specific Cambridge and
COUGHVID tasks \cite{aytekin2024hst}. The DNDT/DNDF paper reported Coswara and
COUGHVID within-dataset AUCs of 0.92 and 0.93, respectively, but its own
cross-dataset table reported Coswara-to-COUGHVID AUC around 0.53
\cite{islam2026robust}. Our conventional external range of 0.523--0.543 is
therefore not evidence that their internal classifiers were fabricated or
incorrect. It is evidence that within-dataset and cross-dataset numbers are not
interchangeable.

The low feature-set overlap complements the drift-adaptive motivation of
Ganitidis \emph{et al.} \cite{ganitidis2025drift}. Our reverse temporal result
also illustrates why AUROC alone is insufficient under a major prevalence or
case-mix change: ranking can remain high while precision, F1, and calibration
are poor. A deployment-oriented analysis must report both discrimination and
the consequences at a stated operating point.

\subsection{Implications for Model Development}

WavLM and CNN--BiGRU did not remove the external gap. Their results show that
the observed source--target difference was not confined to the conventional
feature bank. Architecture comparisons and transportability comparisons remain
separate questions because the cohort, split, threshold, and target population
also determine the estimate.

Fusion should likewise be interpreted as a fitted prediction layer rather than
an automatic benefit of having more modalities. Uniform averaging imposes
equal contribution; validation weighting uses observed validation quality; a
logistic stack learns coefficients from aligned validation predictions. The
test set cannot decide which fusion to report. This boundary is especially
important when participants have missing modalities or when one modality has a
different prevalence and quality distribution.

\subsection{Strengths}

The analysis preserves participant identities, prevents target data from
entering source model selection, distinguishes cough transfer from multimodal
fusion, reports calibration and base-rate-aware operating points, and uses
negative controls. External differences use independent-sample bootstrap
resampling, while DeLong tests are restricted to paired predictions. Results
are accompanied by split, quality, feature, and label-construction audits.

\subsection{Limitations}

The limitations define the scope of inference rather than invalidate the
study. First, Coswara and COUGHVID differ in recruitment, labels, prevalence,
devices, countries, and recording instructions. External failure therefore
cannot be attributed to one isolated mechanism. It is a transportability
result for the complete source-target shift.

Second, COUGHVID is cough-only. We can test whether a Coswara cough branch
transfers, but not whether cough--breath--speech fusion transfers. A second
external dataset with matched modalities and compatible labels would be needed
for that claim.

Third, the temporal protocols are affected by severe changes in label
prevalence and cohort composition. The reverse temporal AUROC is consequently
not a mirror-image estimate of forward deployment. It is retained as an
exploratory stress test with AUPRC and calibration, not presented as a superior
model.

Fourth, the metadata-plus-audio analysis has only 58--61 aligned test
participants across leading candidates. Its intervals are wide. The analysis
can rule out neither a modest benefit nor a modest harm, and it should not be
used as definitive proof that audio adds no information.

Fifth, metadata variables are measured imperfectly and cannot exhaust all
confounding paths. Strong metadata prediction and recording-year importance
establish association with collection context; they do not identify the exact
features used by an audio model. Matching depends on measured covariates and
available overlap and does not reproduce a randomized comparison.

Sixth, COUGHVID labels are imbalanced and may contain label noise. AUROC is
prevalence-invariant in definition but not immune to spectrum and label shift;
AUPRC, calibration, and precision are strongly prevalence-dependent. Both
discrimination and operational metrics are therefore reported.

Seventh, numerous exploratory analyses were performed after the primary model
development. They are labeled exploratory, and their isolated $p$-values are
not interpreted as confirmatory without multiplicity control. The primary
claims rest on the prespecified validation ladder and matched external model
families.

Eighth, the chronological evaluation reused a top-800 bank selected under the
original development split. It prevents later and external labels from entering
feature selection, but it is not a fully nested estimate of an end-to-end model
developed only in the early period. The temporal result is therefore a stress
test of the frozen workflow, not a prospective temporal-validation claim.

Finally, the study is retrospective and uses public crowdsourced audio. It does
not evaluate prospective workflow integration, patient outcomes, fairness in a
deployment population, or comparison with an approved clinical test. Decision
curves describe the present sample and utility assumptions; they do not show
that deployment would cause or prevent clinical harm.
